\chapter{Repository and Module Map}
\label{ch:repo-map}

This appendix records the exact repository and physical-module boundary used for this edition.
It is intentionally narrower than an organization-wide project directory: every listed sibling
has a pinned revision and a defined role in dependency, verification, examples, or consumption.
All sibling trees were read-only during the book audit.

The revisions below are immutable edition pins, not a promise that every sibling's branch head
will remain there.  For example, the local \texttt{cna-template} \texttt{next} branch advanced to
\texttt{4d0a2c8a} after this audit; this edition continues to cite and reason from
\texttt{1d86681}, which remains available in that repository.

\section{Pinned repository set}

\begin{small}
\begin{longtable}{>{\raggedright\arraybackslash}p{2.25cm}>{\ttfamily}p{1.45cm}>{\raggedright\arraybackslash}p{7.7cm}}
\toprule
\textbf{Repository} & \textbf{Pin} & \textbf{Role in this edition} \\
\midrule
\endfirsthead
\toprule
\textbf{Repository} & \textbf{Pin} & \textbf{Role in this edition} \\
\midrule
\endhead
\texttt{cna} & 7a64362e & Primary C\texttt{++}23 XNA-style framework; source, build,
  renderer, platform and verification subject of the book. \\
\texttt{sharp-runtime} & f827a6c5 & .NET-shaped C\texttt{++} foundation consumed by CNA;
  object/type, collections, IO, threading/task and compatibility evidence. \\
\texttt{easy-gl} & 0b46d35 & Shared implementation beneath five public EasyGL renderer
  identities; owns portable GL objects and draw machinery. \\
\texttt{meta-gl} & 571d3a6 & Lower-level typed OpenGL wrapper consumed by easy-gl; not a
  separate CNA renderer identity. \\
\texttt{free-direct} & 934f72f & SDL3-based DirectDraw/DirectSound/DirectPlay compatibility
  layer used by CNA's \texttt{FREEDIRECT} renderer and the Blupi case study. \\
\texttt{free-api} & 53d7a31 & Minimal Win32-style support layer consumed by free-direct; a
  second-level dependency, not a direct CNA module. \\
\texttt{xna4-spec} & 8f61207 & Converted XNA documentation corpus: 544 indexed XML type files
  across 19 documented namespaces; useful inventory, not a live CNA gate. \\
\texttt{cna-samples} & 149a19f & Official-sample translation catalog: 63 completed ports, 23
  tracked placeholders, and 67 explicitly ignored source-archive entries at the pin. \\
\texttt{cna-examples} & fa0b49a & One CNA-native catalog application with 13 areas, 79
  categories, and 249 mechanically checked demo screens. \\
\texttt{cna-template} & 1d86681 & Minimal consumer skeleton; evidence for how a new application
  composes the framework rather than for framework correctness. \\
\texttt{cna-craft} & 7d7f186 & Larger game consumer with desktop/Web pressure; useful
  integration evidence, not part of CNA's own test suite. \\
\texttt{cna-extended} & 2ff3cff & Extension-library consumer modeled after
  MonoGame.Extended; outside the core API but relevant to package/composition boundaries. \\
\bottomrule
\end{longtable}
\end{small}

The dependency arrows are not all equivalent:

\begin{verbatim}
cna ---> sharp-runtime
 |----> easy-gl ---> meta-gl          (selected GL family)
 `----> free-direct ---> free-api     (selected FREEDIRECT family)

xna4-spec, cna-samples, cna-examples          evidence/catalog peers
cna-template, cna-craft, cna-extended         downstream consumers
\end{verbatim}

Renderer-specific third parties fetched or discovered by CNA are not promoted to sibling
repositories in this diagram.  They remain dependencies of their owning renderer module, with
their platform and acquisition gates recorded in Part~\ref{part:renderers}.

\section{Physical CNA framework modules}

The root \path{modules/CMakeLists.txt} composes fourteen framework/extension modules plus the
renderer tree.  Directory location is source ownership: production code belongs under the
owner's \texttt{src/}, public and public-internal contracts under its \texttt{include/}, and its
tests/examples alongside them.  A configure-time validator rejects translation units outside a
declared module tree.

\begin{center}
\begin{tabular}{p{3.1cm}p{8cm}}
\toprule
\textbf{Physical module} & \textbf{Primary ownership} \\
\midrule
\texttt{core} & CNA helper/marker surface and root framework primitives \\
\texttt{math} & vectors, matrices, geometry, curves and packed math values \\
\texttt{graphics} & device facade, resources, effects, model runtime and shared renderer contract \\
\texttt{input} & keyboard, mouse, gamepad, touch and SDL input extensions \\
\texttt{audio} & SoundEffect/XACT-facing audio objects and SDL3\_mixer integration \\
\texttt{media} & Song/MediaPlayer/Video and FFmpeg-conditioned paths \\
\texttt{content} & ContentManager, XNB, CNJ, loose readers and glTF import core \\
\texttt{storage} & StorageDevice/StorageContainer and local-filesystem behavior \\
\texttt{runtime} & Game, GameWindow, components, services and timing/lifecycle composition \\
\texttt{devices} & Microsoft/XNA-shaped device and sensor surface \\
\texttt{devices-ext} & CNA-native camera/dialog/system/clipboard-style host services \\
\texttt{graphics-ext} & CNA graphics extensions, including the former extension-layer library \\
\texttt{gamer-services} & local gamer identity, guide, achievements and leaderboards; conditional
  on networking enablement \\
\texttt{net} & sessions, gamers, packets and ENet-backed delivery/discovery; conditional \\
\bottomrule
\end{tabular}
\end{center}

The interface target \texttt{CNA} composes the framework modules, extension umbrellas, shared
build flags, sharp-runtime closure, and exactly one selected renderer target.  Individual module
aliases support narrower link-closure probes.  \texttt{CNA::CnaExt} composes
\texttt{graphics-ext} and \texttt{devices-ext}; it is distinct from the empty \texttt{CNAEXT}
declaration marker and from the \texttt{CNA\_CNAEXT} build option that enables selected extension
products.

\section{Renderer module tree}

\path{modules/renderers/} contains one common-contract area plus 42 implementation families.
Only EasyGL maps multiple public identities to one family.

\begin{small}
\begin{longtable}{>{\raggedright\arraybackslash}p{3.2cm}>{\raggedright\arraybackslash}p{8.2cm}}
\toprule
\textbf{Cluster} & \textbf{Physical family directories} \\
\midrule
\endfirsthead
\toprule
\textbf{Cluster} & \textbf{Physical family directories} \\
\midrule
\endhead
OpenGL & \texttt{easygl}, \texttt{opengles1}, \texttt{opengl1}, \texttt{opengl2},
  \texttt{opengl4}, \texttt{portablegl} \\
Native modern & \texttt{vulkan}, \texttt{webgpu}, \texttt{sdl-gpu}, \texttt{metal} \\
Abstraction layers & \texttt{bgfx}, \texttt{magnum}, \texttt{llgl}, \texttt{diligent},
  \texttt{sokol}, \texttt{wicked}, \texttt{fna3d} \\
DirectX / retro & \texttt{directx1}, \texttt{directx2}, \texttt{directx3},
  \texttt{directx5}, \texttt{directx6}, \texttt{directx7}, \texttt{directx8},
  \texttt{directx9}, \texttt{directx10}, \texttt{directx11}, \texttt{directx12},
  \texttt{glide}, \texttt{freedirect} \\
2D / vector / Windows & \texttt{sdl-renderer}, \texttt{skia}, \texttt{blend2d},
  \texttt{openvg}, \texttt{direct2d}, \texttt{gdi} \\
Web DOM & \texttt{canvas}, \texttt{html-dom}, \texttt{svg-dom} \\
Diagnostic / CPU & \texttt{headless}, \texttt{software}, \texttt{stub} \\
\bottomrule
\end{longtable}
\end{small}

The physical tree alone does not establish a public selector, platform admissibility, resource
factory, or semantic parity.  Appendix~\ref{ch:feature-matrix} joins directory families to all
46 public identities and their narrow implementation behavior.

\section{How to extend this map without making it stale}

For a later edition, pin every repository before counting it; derive module and renderer sets
from CMake/source rather than an organization README; distinguish direct dependency, transitive
dependency, evidence peer and downstream consumer; and record the command/query behind every
count.  A local worktree of \texttt{cna} is not a sibling project, and an organization repository
with no dependency or evidence edge is outside this bounded map.
